\setcounter{chapter}{13} % "N" (14th letter, but latex first increments -> 14-1=13)
\chapter{Notation \& Conventions} \label{ch:notation}
\markboth{Notation \& Conventions}{Notation \& Conventions}

Throughout the thesis, uppercase Latin letters such as \(A\), \(S\), and \(Q\)
denote matrices. At the same time, they are often interpreted as concrete
representations of linear maps between finite-dimensional inner-product spaces.
Vectors are denoted by lowercase letters such as \(x\) and \(b\). For a least-
squares problem we usually write \(A \in \R^{n \times d}\) with \(n \ge d\), and we
consider the overdetermined problem
\[
	\min_{x \in \R^d} \|Ax-b\|_2.
\]

The symbol \(S\) is reserved for a sketching matrix, or equivalently for a linear
map that compresses the ambient data space. A sketched least-squares problem
therefore has the form
\[
	\min_{x \in \R^d} \|SAx-Sb\|_2.
\]

Unless stated otherwise, \(\|\cdot\|_2\) denotes the Euclidean norm for vectors and
the induced operator norm for matrices, while \(\|\cdot\|_F\) denotes the Frobenius
norm. The condition number of a nonsingular matrix \(M\) in the Euclidean norm is
written as
\[
	\kappa_2(M) = \|M\|_2 \|M^{-1}\|_2
\]
and is interpreted as the sensitivity of the associated linear system to
perturbations.

Floating-point quantities are discussed using the standard model of rounding to
nearest. The unit roundoff of a format is denoted by \(u\). When mixed precision
is used, different unit roundoffs may appear in different parts of the same
algorithm. These are distinguished by subscripts such as \(u_\ell\) for lower
precision and \(u_h\) for higher precision.

In numerical experiments, the terms \enquote{solution error} and \enquote{residual error} are used in the following sense:
\begin{itemize}
	\item solution error compares an approximation \(\hat x\) directly with a reference solution \(x_\star\), typically through \(\|\hat x - x_\star\|_2\) or a relative variant,
	\item residual error measures how well \(\hat x\) approximates the data vector \(b\) through quantities such as \(\|A\hat x - b\|_2\).
\end{itemize}

Gaussian sketches, structured sketches, and mixed-precision variants are defined
more carefully in the later chapters. The notation fixed here remains standard
throughout the text. We also move freely between matrix language and the
corresponding operator viewpoint whenever this makes the numerical question
clearer.
